\section{What the Second Hop Costs}
\label{sec:ch07-what-the-second-hop-costs}
\index{latency!federated query|(}%

Everyone who meets federation asks the same question, and it is a fair one: what
does all this cost? The sample can answer it, in the narrow way a sample can.

I measured four things on one machine in one sitting: five warm-up requests then
fifty timed ones for each row, timed with \code{curl}'s own
\code{\%\{time\_total\}}. Medians, in milliseconds.

\begin{center}
\begin{tabular}{lrrrr}
\toprule
 & min & median & p90 & max \\
\midrule
Router, both subgraphs   & 4.14 & 4.69 & 5.84 & 24.14 \\
Router, catalog only     & 2.87 & 3.27 & 3.82 & 20.77 \\
Catalog, direct          & 0.70 & 0.81 & 0.96 &  1.43 \\
Reviews \code{\_entities}, direct & 0.78 & 0.92 & 1.07 & 21.12 \\
\bottomrule
\end{tabular}
\end{center}

Three of the four maxima are around 20 milliseconds and they are the machine,
not the graph: each of those rows has exactly one such outlier and nothing else
near it. Read the medians.

\index{latency!second subgraph hop}%
\medskip\noindent\textbf{The second hop costs about 1.4 milliseconds.} That is
4.69 against 3.27, the same router answering the same root field with and
without a selection that crosses a boundary. It lines up with what the Reviews
subgraph costs when I ask it directly, 0.92, which is the result you would
predict: the second fetch is a serial dependency, so its cost adds rather than
hides behind anything.

\index{latency!router overhead}%
\medskip\noindent\textbf{The router costs more than the second subgraph does.}
3.27 against 0.81 for the same single-subgraph answer. Roughly two and a half
milliseconds of parsing, normalisation, planning, HTTP client work and container
networking, against one and a half for an entire extra service.

\medskip
I did not expect the second one and I have stopped being surprised by it. The
mental model most people arrive with says a federated query is expensive because
it fans out, so the thing to optimise is the number of subgraphs a query
touches. On these numbers that instinct would send you after the smaller half of
the cost.

\index{latency!measurement caveats}%
Three things bound how far those numbers travel. The router runs in a container
and reaches the subgraphs across Docker's bridge to the host, while the direct
rows are two processes talking on one host, so the router's 2.5 milliseconds
includes a network hop the subgraphs do not pay and is an upper bound rather
than a measurement of its own work. The subgraphs answer out of dictionaries
holding three products, with no database underneath and nothing large to
serialise. And the machine is mine, which
chapter~\ref{ch:data-without-the-n-plus-1} already established is worth saying
out loud, because it measured the same code at ten milliseconds one afternoon
and twenty the next purely because something else was running.

What survives to your hardware is the ratio and its direction, not the
milliseconds. A cross-boundary hop against two tiny in-memory subgraphs is
cheap; the fixed cost of having a router at all is not, and it is paid by every
query including the ones that never leave one subgraph.

\index{query plan!as a cost estimate}%
That has a practical consequence worth carrying into
chapter~\ref{ch:enter-the-router}. The plan of
section~\ref{sec:ch07-the-plan} tells you the number of hops before you deploy
anything, and counting them costs nothing. It will not tell you what any of
them cost. A plan with two
fetches where one of them queries a table nobody indexed is a fast plan over a
slow query, and the plan cannot see the difference.
Chapter~\ref{ch:resilience-and-performance} is where the two are measured
together.

\index{latency!federated query|)}%
